\clearpage\par\bigskip\Needspace{6\baselineskip}
\pdfbookmark[0]{Remarks on the Analysis Residuorum}{bemerkungen-analysis-residuorum}
\begin{center}\large\bfseries REMARKS ON THE ANALYSIS RESIDUORUM.\end{center}
\BookRule{.09\linewidth}
\renewcommand{\BookRunningTitle}{Historical Editor's Remarks on the Analysis Residuorum}

The two preceding treatises are taken from an extensive manuscript entitled Analysis Residuorum, which probably dates from 1797 or 1798; a complete revision of it later gave rise to the Disquisitiones Arithmeticae. The full title of Chapter Six reads:

Solution of the congruence $x^m-1\equiv0$ and the equation $x^m-1=0$; with elucidations on the theory of regular polygons.

The essential content of its second part (\textsection\textsection.\ 253--278) passed into the seventh section of the Disqq. Arithm.

Chapter Seven is also partly preserved: Various Applications of Certain Preceding Investigations (\textsection\textsection.\ 279--302). It is divided into the following subsections:

On Transformations of Common Fractions (\textsection\textsection.\ 279--281).

On the Conversion of Common Fractions into Decimals (\textsection\textsection.\ 282--292).

On the Solution of the Indeterminate Equation $xx=a+by$ (\textsection\textsection.\ 293--297).

On the Solution of the Indeterminate Equation $axx+byy=c$ (\textsection\textsection.\ 298--301).

On Finding the Divisors of Numbers (\textsection.\ 302; the following sheets are missing).

All this was incorporated almost word for word into the sixth section of the Disqq. Arithm.

The two sections given here treat the subjects which, as appears from the preface and Articles 11, 44, 61, 62, 65, 84 of the Disqq. Arithm., were intended to form the eighth section of that work. It should, however, be noted that this plan was later changed again: among the manuscripts is a fragment headed Section Eight: A Fuller Consideration of Certain Investigations Concerning the Division of the Circle. This fragment
\SourceLeaf{251}
begins with Art.\ 367 and was therefore intended to continue the Disqq.\ Arithm.; but the content of the few articles still preserved later passed into the treatise Summatio quarundam serierum singularium, and for this reason the fragment is excluded from the present edition.

In the preceding printing of the two parts of the Analysis Residuorum, the original text has essentially been retained faithfully, although its form cannot be called ready for printing; the following remarks indicate the most important alterations and add some elucidations.

\textsection.\ 237.\ Compare Disqq.\ Arithm.\ artt.\ 61, 62.

\textsection.\ 239.\ Compare Disqq.\ Arithm.\ artt.\ 53, 54, 65.

\textsection.\ 241.\ If $n=2^\nu$ and $\nu\geqq3$, there is indeed no number $\rho$ of the stated kind, but this does not substantially change the investigation as a whole.

\textsection.\ 251.\ The difficulty noted here was presumably to be removed by introducing higher powers of $p$ as moduli. Compare \textsection\textsection.\ 363, 372, 373.

\textsection.\ 332.\ The assumption that the modulus is a prime number is retained through \textsection.\ 372 inclusive.

\textsection.\ 338.\ The incomplete citation may refer to Disqq.\ Arithm.\ art.\ 44.

\textsection\textsection.\ 344--346.\ Of the two proofs present in the manuscript, the first, introduced by the words ``we now set about proving'' and based on a closer examination of the expressions $(1^\alpha2^\beta3^\gamma\ldots)$, has here been omitted entirely in accordance with the author's own instruction (``The whole preceding proof, together with another proof of the preceding theorem which we intended to add, is to be suppressed, because we have discovered another that is infinitely simpler. It rests on this foundation''); moreover, in the printing above the second proof has been shortened by considering the expansion of $\frac{x\,\mathrm{d}P}{P\,\mathrm{d}x}$ instead of that of $\frac{x\,\mathrm{d}P}{\mathrm{d}x}$, thus also avoiding a reference in the original to the suppressed first proof.

\textsection.\ 348.\ The expression radix prima is to be understood here in the same sense as radix propria in the treatise Summatio quarundam serierum singularum, art.\ 11. --- The assertion that the coefficients $A'.B'\ldots$ of the expanded product are rational integers refers to Chapter Six, where, however, the theory of the equation $x^\tau-1=0$ is treated only for the case where $\tau$ is prime; the form of the proof in \textsection.\ 349 first leads to the following supplement. If the expanded product is put into the form (valid for all roots of the equation $\theta^\tau=1$)
\[
S=E+F\theta+\ldots+N\theta^{\tau-1}
\]
then the coefficients $E,F\ldots N$ are polynomials in $x$ with rational integer coefficients. Further, since the product remains unchanged when $\theta$ is replaced by $\theta^k$, where $k$ is any number coprime to $\tau$, the same holds for the expression $S$. It follows without difficulty that all powers occurring in $S$ of $\theta$ whose exponents $s$ have one and the same greatest common divisor with $\tau$ must also have identical coefficients. Finally, since every sum of such powers $\theta^s$ is always an integer, it is evident that the expression $S$, and hence also the product in question, is a polynomial in $x$ with integer coefficients, as was to be shown. This consideration also establishes the correctness of the remark at the end of the paragraph. Other reasons nevertheless suggest that the author already knew at that time the general theorem on the transformation of symmetric functions (Demonstratio nova altera theorematis omnem functionem etc.\ art.\ 4), from which the propositions above follow immediately.

\textsection.\ 352.\ The notation $R\equiv S(\mathrm{mod.}\,P)$, or also $R\equiv S(\mathrm{mod.}\,P,p)$, means here and in what
\SourceLeaf{252}
follows that the difference $R-S$ has, modulo $p$, the divisor $P$. --- The incomplete citation may refer to Disqq.\ Arithm.\ art.\ 49.

\Needspace{10\baselineskip}\textsection.\ 354.\ Multiplication by $x^\nu-1$ shows that the sums of equal powers of the roots of the two equations $(P,\rho^{k\nu+t})=0$, $(P,\rho^t)=0$ are congruent to one another $(\mathrm{mod.}\,p)$, and from this follows the congruence $(P,\rho^{k\nu+t})\equiv(P,\rho^t)(\mathrm{mod.}\,p)$ as soon as $m<p$ (compare \textsection.\ 244). But if $m\geqq p$, the coefficient of the power $x^{m-p}$ in an equation can no longer be determined modulo $p$ from the given power sums of its roots, because it carries the factor $p$ in the formulas of \textsc{Newton} used for this purpose. Indeed, from the congruence of each pair of equal power sums of the roots of the equations $A=0$, $B=0$, one may in general infer only that $A\equiv\mathfrak A^p\mathfrak C$, $B=\mathfrak B^p\mathfrak C(\mathrm{mod.}\,p)$, where $\mathfrak C$ denotes the greatest common divisor of the two functions $A,B$ modulo the prime number $p$, while $\mathfrak A$ and $\mathfrak B$ are entirely undetermined functions. Presumably the author derived the general validity of the proposition from the theory of the transformation of symmetric functions, and specifically from the following proposition: if, with respect to any modulus $p$, the difference $R(x)-S(x)$ is divisible by the function $P(x)$, and $a,b,c\ldots$ are the roots of the equation $P(x)=0$, then the functions
\[
\begin{gathered}
(x-R(a))(x-R(b))(x-R(c))\ldots\quad\text{and}\\
(x-S(a))(x-S(b))(x-S(c))\ldots
\end{gathered}
\]
are congruent to one another modulo $p$.

\textsection.\ 355.\ It is shown in \textsection.\ 368 that $P$ and $\frac{\mathrm{d}P}{\mathrm{d}x}$ have no common divisor if $P$ contains no factor more than once.

\textsection\textsection.\ 358, 359.\ The note placed below the text is taken from a separate sheet, which was probably intended to replace \textsection.\ 359, already crossed out in the manuscript.

\textsection.\ 360.\ An inaccuracy in the manuscript has been corrected in the statement of the theorem.

\textsection.\ 361.\ Here the exponent $\frac1k$ in the notation $(\xi,\rho^{\frac1k})$ means any positive integer $k'$ such that $kk'\equiv1(\mathrm{mod.}\,\nu)$, where $\nu$ is the least positive integer for which $x^\nu-1$ is divisible by $\xi$ modulo $p$; it must be assumed here that $\xi$ is not divisible by $x$ modulo $p$, and also that $k$ is coprime to $\nu$. The correctness of the assertion that $\xi'$ is divisible by $(\xi,\rho^{\frac1k})$ $(\mathrm{mod.}\,p)$ follows from \textsection.\ 354.

\textsection.\ 363.\ The concluding remark presumably concerns the introduction of moduli that are powers of the prime number $p$; compare \textsection\textsection.\ 251, 372, 373.

\Needspace{14\baselineskip}\textsection.\ 367.\ The roots of the equation $x^3+xx-2x-1=0$ are the periods of two terms into which the roots of the equation $\frac{x^7-1}{x-1}=0$ divide. The same example also occurs on a separate sheet, where the principal result of \textsection\textsection.\ 362, 363 is stated under the title ``the golden theorem''.

\textsection.\ 371.\ This paragraph was intended to contain an example, but it was not worked out.
\begin{flushright}R.\ \textsc{Dedekind}.\end{flushright}
\BookRule{.28\linewidth}
